\documentclass[11pt]{article}

% --- Packages ---
\usepackage[utf8]{inputenc}
\usepackage[T1]{fontenc}
\usepackage{geometry}
\geometry{letterpaper, margin=1in}
\usepackage{setspace}
\doublespacing
\usepackage{graphicx}
\usepackage{float}
\usepackage{amsmath}
\usepackage{natbib}
\bibliographystyle{unsrtnat}
\usepackage{hyperref}
\hypersetup{colorlinks=true, linkcolor=blue, citecolor=blue, urlcolor=blue}
\usepackage{booktabs}
\usepackage{siunitx}
\usepackage{xcolor}
\graphicspath{{figures/}}

% --- Custom commands ---
\newcommand{\msix}{m$^6$A}

\newcommand{\LONE}{LINE-1}

% --- Title ---
\title{%
 Poly(A) tail retention of LINE-1 transcripts under arsenite stress \\
 covaries with m$^6$A density and structural completeness
}
\author{%
 Junsoo Kim\textsuperscript{1,*},
 Sojeong Lee\textsuperscript{1,2},
 Seungbeom Han\textsuperscript{1,3}
 \\[6pt]
 \textsuperscript{1}Center for RNA Research,
 Institute for Basic Science (IBS), \\
 Seoul National University, Seoul, Republic of Korea \\
 \textsuperscript{2}School of Biological Sciences,
 Seoul National University, \\
 Seoul, Republic of Korea \\
 \textsuperscript{3}Interdisciplinary Program in Bioinformatics, \\
 Seoul National University, Seoul, Republic of Korea \\[3pt]
 \textsuperscript{*}Corresponding author: junsoopablo@snu.ac.kr
}
\date{}

% ============================================================
\begin{document}
\maketitle

% ============================================================
% ABSTRACT
% ============================================================
\begin{abstract}

\LONE{} (L1) retrotransposons comprise 17\% of the human genome and are
pervasively transcribed across human cell types, yet how RNA
modifications shape L1 transcript fate under cellular stress remains
unknown. Here, we use Oxford Nanopore direct RNA sequencing across nine
human cell lines (23 libraries) to simultaneously measure \msix{},
poly(A) tail length, and transcript integrity on individual L1 molecules.
L1 transcripts carry higher \msix{} density than non-L1 transcripts,
beyond sequence-composition expectation, and \msix{} density varies
with evolutionary age. Under arsenite stress, L1 poly(A) tails shorten
while non-L1 transcripts remain unchanged. Young, full-length L1
retain their poly(A) tails, and overlapping sequence features---structural
completeness, EN domain retention, high \msix{} density, and long
consensus span---each independently confer this protection when present
in ancient elements. A composite score integrating these features
separates stressed ancient L1 into subgroups spanning 66--137~nt in
median poly(A). \msix{} density covaries with poly(A) retention
specifically under stress, with low-\msix{} reads more frequently
reaching critically short poly(A) lengths. Together, these findings
reveal that \msix{} density, uncorrelated with poly(A) under basal
conditions, becomes a stress-specific correlate of poly(A) retention
among structurally incomplete L1.

\end{abstract}

% ============================================================
% INTRODUCTION
% ============================================================
\section{Introduction}

Long interspersed nuclear element 1 (\LONE{}, L1) retrotransposons are the
most abundant autonomous mobile elements in the human genome, constituting
approximately 17\% of total genomic sequence \citep{Bourque2018}. Although
most of the $\sim$500,000 L1 copies are ancient, truncated remnants
incapable of retrotransposition, a small number of evolutionary young,
full-length elements (L1HS and L1PA subfamilies) retain
retrotransposition competence \citep{Lanciano2020}. Beyond parasitic
replication, L1 transcripts serve regulatory functions, including nuclear
scaffolds in embryonic stem cells \citep{Percharde2018}, long-range
enhancers \citep{LiX2024}, and contributors to long noncoding RNA
biogenesis \citep{Kapusta2013}. Notably, ancient L1 fragments are
preferentially exapted as cis-regulatory elements: transposable
element--derived sequences contribute $\sim$25\% of candidate
cis-regulatory elements in the human genome, with LINE families among
the largest contributors \citep{Du2024encode}. Enhancer-associated
copies retain higher sequence conservation and longer intact length
than non-functional counterparts \citep{Nishihara2019}, and the
probability of enhancer exaptation increases with element age
\citep{Simonti2017}. Thus, L1 RNA is both a potentially
deleterious mobile element and a functional regulatory molecule. Multiple
layers of host regulation---DNA methylation, histone modifications, and
RNA-based mechanisms---act to restrain L1 activity \citep{Deniz2019}.

$N^6$-methyladenosine (\msix{}) is the most prevalent internal
modification of mammalian mRNA and a key regulator of L1 biology. \msix{}
has dual roles on L1 retrotransposons. METTL3-deposited \msix{} in the L1
$5'$ UTR promotes cap-independent translation of ORF1p and ORF2p via eIF3
recruitment, enhancing retrotransposition \citep{Hwang2021}. Hwang
et al. further showed that \msix{} consensus motifs in the L1 $5'$ UTR
have been under positive evolutionary selection, indicating that L1
co-opts the host \msix{} machinery. Conversely, the host cell exploits
the same marks to suppress L1 activity---an \msix{}-based innate defense
mechanism that also targets other transposable elements
\citep{Chelmicki2021}---through multiple readers. The nuclear reader
YTHDC1 recognizes \msix{}-marked L1 RNA and recruits SETDB1 for H3K9me3
chromatin silencing \citep{Liu2021}, while the cytoplasmic reader YTHDF2
directs \msix{}-L1 RNA to autophagy-mediated decay \citep{Tian2024}.
Beyond this antagonism, \msix{} on intronic L1 participates in gene
regulation, as SAFB engages \msix{}-marked intronic L1 to modulate long
gene expression \citep{Xiong2021}. This range of \msix{}-dependent
functions has been the prevailing framework for understanding \msix{} on
L1.

However, several questions remain unaddressed. First, most studies focus
on young L1 elements (L1HS, L1PA1--2). The \msix{} landscape of
\emph{ancient} L1---approximately 86\% of all L1 copies and the
predominant source of exapted regulatory elements---remains largely
unexplored, including whether the DRACH motif context required for
METTL3 recognition is preserved despite millions of years of sequence
divergence. Second, prior \msix{} quantification of L1 has relied on
short-read approaches (MeRIP-seq, miCLIP) with severe mapping ambiguity
across the repetitive L1 family \citep{Lanciano2020}; single-molecule
\msix{} quantification on L1 has not been achieved. Third, L1 \msix{}
behaviour under cellular stress and its relationship to L1 RNA stability
have not been investigated, despite evidence that arsenite stress drives
L1 ORF1p accumulation in stress granules \citep{Goodier2007} and that
\msix{} mediates stress-granule triage of mRNAs \citep{Anders2018}.

Arsenite stress provides a useful context for studying L1 RNA fate.
Sodium arsenite triggers rapid stress granule formation and activates a
non-canonical RNA decay pathway, XRN1-mediated $5'$ exonucleolytic
digestion that leaves poly(A) tails intact \citep{deMuro2024}. Canonical
deadenylation is suppressed under arsenite, as the CCR4-NOT complex is
sequestered into stress granules \citep[preprint]{Bhatt2025} and the
deadenylase adaptors Tob and Pan3 are degraded \citep{Funakoshi2015}.
Whether non-translated L1 fragments---intronless, AU-rich RNA
molecules that lack the ribosomal protection of young, actively
translated elements---are subject to distinct decay kinetics under these
conditions has not been examined.

Here, we leverage Oxford Nanopore direct RNA sequencing (DRS) across nine
human cell lines to simultaneously measure \msix{}, poly(A) tail length,
and transcript integrity on individual L1 molecules. We report that L1
carries elevated \msix{} density across evolutionary age, that arsenite
selectively shortens L1 poly(A) tails while young L1 retain theirs, and
that overlapping sequence features---including \msix{} density---distinguish
protected from vulnerable elements, revealing a stress-specific dimension
of \msix{} on L1 distinct from its established roles in retrotransposition
and silencing.

% ============================================================
% RESULTS
% ============================================================
\section{Results}

% ------------------------------------------------------------
\subsection{L1 RNA carries elevated \msix{} density across human cell lines}

We applied Oxford Nanopore direct RNA sequencing (DRS) to nine human cell
lines and identified 54,234 L1 reads across 12,125 loci using a
two-stage filter that removes $\sim$90\% of host gene read-through (see
Methods). Of these, 38,544 were non-embedded and 15,690 were embedded
within host gene exons (Table~S1). Ancient elements (L1MC,
L1ME, L1M) comprised 93\% and young elements (L1HS, L1PA1--3) 7\% of the
non-embedded reads (Supplementary Fig.~S1). Long-read alignment largely
resolves the multi-mapping ambiguity that confounds short-read L1
studies \citep{Lanciano2020}: ancient L1 reads had a median mapping
quality (MAPQ) of 60 with $>$80\% uniquely mapped (Supplementary
Table~S1). Young L1 reads, owing to $>$97\% inter-element sequence
identity, had lower MAPQ (median 9--29 for L1HS--L1PA2), but the
kilobase-scale read lengths of DRS still anchor most reads to individual
loci---an advantage unavailable to 150-bp short reads. Because our
primary conclusions rely on subfamily-aggregated metrics (\msix{}/kb,
poly(A) length) rather than locus-level quantification, the residual
multi-mapping in young L1 does not affect the reported enrichments.
Among non-embedded L1, 61\% resided within introns and 39\% in
intergenic regions, with similar mapping quality and poly(A)
distributions (Supplementary Table~S1).

L1 transcripts carried \msix{} at a density of 2.89 sites/kb, compared
with 1.62 sites/kb for non-L1 transcripts (randomly subsampled reads
from the same libraries, predominantly mRNA; see Methods),
a 1.79-fold enrichment (Fig.~1a). This exceeds sequence-composition
expectation, as L1 has only 1.22-fold more DRACH motifs per kb than
non-L1 transcripts. The enrichment was consistent across
all nine cell lines and persisted at every probability threshold tested
, with high-confidence \msix{} calls occurring
more frequently on L1. \msix{} density also correlated with evolutionary age.
Young L1 (L1HS, L1PA1--3) exhibited the highest density (4.16 sites/kb),
followed by ancient L1 (2.69 sites/kb) and non-L1 transcripts (1.62
sites/kb, Fig.~1a). The Young${>}$Ancient
gradient and L1/non-L1 enrichment were consistent across all nine cell
lines (all $P < 10^{-7}$, Supplementary Fig.~S4a), and subfamily-level
\msix{}/kb rankings were highly concordant between cell line pairs
(median Spearman $\rho = 0.85$, Supplementary Fig.~S4b), indicating
sequence-intrinsic features. \msix{} density showed modest positional
variation along the L1 consensus, with ORF1-derived fragments carrying
the highest density.

Orthogonal validation with RNA004 chemistry and the Dorado basecaller
reproduced the L1 enrichment. In length-matched comparisons, L1 reads
carried 1.5--1.8-fold higher \msix{}/kb than non-L1 reads (Supplementary Fig.~S6),
consistent with RNA002 results, and the Young${>}$Ancient gradient was
recapitulated (1.81-fold). L1 \msix{} showed strong DRACH-motif
specificity, with the DRACH motif hierarchy nearly identical between L1
and non-L1 (Spearman $\rho = 0.91$, $P = 1.7 \times 10^{-7}$), with
GGACT ranking first in both. Among 539 position-reproducible \msix{} sites
on L1 ($\geq$50\% of reads methylated at a given genomic position), 93\%
occurred at DRACH motifs (Supplementary Fig.~S6).

To evaluate ancient L1, we asked whether the positional \msix{} landscape
mirrors that of young elements. When DRACH sites were mapped to the L1
consensus sequence, per-position methylation rates of ancient and young
L1 were correlated (Spearman $\rho = 0.44$, $P = 1.3 \times 10^{-6}$,
$n = 112$ consensus bins): positions that are \msix{} hotspots
in young L1 showed 1.78-fold higher methylation in ancient copies than
non-hotspot positions ($P = 5.0 \times 10^{-94}$). Methylation rates at
individual DRACH sites were predicted by flanking sequence conservation.
Ancient L1 DRACH sites retaining Young-L1-like flanking sequences
($\leq$2 mismatches in an 11-mer window) were methylated at twice the
rate of diverged sites (9.95\% vs.\ 4.98\%; $P = 1.3 \times 10^{-96}$),
in line with reports that flanking sequence and local RNA
structure influence methylation \citep{Luo2022,Shachar2024}. The
Young${>}$Ancient gradient arises from two multiplicative factors. In
length-matched reads (2--5~kb), young L1 carried 1.41-fold more DRACH
motifs per kb and showed 1.64-fold higher per-DRACH methylation (Fisher
$P < 10^{-30}$ for both). Together, these results indicate that L1
retrotransposons carry \msix{} enrichment that exceeds
sequence-compositional expectation, shows an evolutionary gradient,
and displays consistent positional patterns across ancient elements.

% --- Figure 1 ---
\begin{figure}[H]
\centering
\includegraphics[width=\textwidth]{fig1.pdf}
\caption{\textbf{L1 \msix{} enrichment and arsenite-selective poly(A) shortening.}
See Figure Legends for full description.}
\label{fig:enrichment_shortening}
\end{figure}

% ============================================================
% RESULTS Part 2
% ============================================================
\subsection{Arsenite triggers selective poly(A) shortening of non-embedded L1 transcripts}

Having established that L1 carries elevated \msix{}, we next asked how
L1 transcripts behave under cellular stress. We analysed ONT DRS data
from HeLa cells treated with 0.5~mM sodium arsenite for 60 minutes
\citep{deMuro2024} and compared poly(A) tail lengths between stressed
and unstressed conditions. L1 poly(A) tails shortened by a median of
31~nt (122.5 to 91.5~nt, Mann--Whitney $P = 2.6 \times 10^{-17}$;
Supplementary Fig.~S1a),
whereas non-L1 transcripts showed no significant change
($\Delta = +2.5$~nt, $P > 0.05$, Fig.~1b). The L1-specific effect
($\Delta\Delta = -33.5$~nt) indicates that shortening is restricted to
L1 rather than a global change. To test whether this reflects a
compositional shift toward loci with inherently shorter tails, we
tracked 2,268 L1 loci detected exclusively under arsenite across nine
unstressed cell lines. These loci had a median poly(A) of 77.5~nt under
arsenite but 125.5~nt in other cell lines ($\Delta = -48.0$~nt,
$P = 2.2 \times 10^{-35}$, Supplementary Fig.~S1a), confirming a stress-dependent
effect rather than locus-intrinsic properties. The shortening was
comparable in intronic and intergenic L1 ($\Delta = -32.9$ and
$-27.1$~nt, respectively; Supplementary Table~S1), persisted at every
read-length stratum, and was absent in annotated lncRNAs processed
through identical pipelines ($\Delta = -2.9$~nt, Cohen's $d = -0.06$;
Supplementary Fig.~S7a--b), establishing it as L1-specific rather than
a general feature of non-coding RNA.

Shortening was confined to ancient L1 elements (L1MC, L1ME, L1M,
$\Delta = -33$~nt, $P = 1.7 \times 10^{-18}$), whereas young L1 were
unaffected ($\Delta = -0.1$~nt, $P = 0.77$, Fig.~1c), even after
matching for \msix{}/kb. Within ancient subfamilies, vulnerability
followed an evolutionary gradient: deeply ancient elements (L1M5,
$\Delta = -49$~nt; L1ME4a, $-37$~nt; L1MC3, $-64$~nt) were most
affected, whereas intermediate-age subfamilies (L1PA5--7) were
unaffected ($P > 0.1$). Structural completeness
was a strong determinant: full-length ancient elements were resistant
($\Delta = -9$~nt, ns; $n = 114$), whereas partial elements showed
significant shortening ($\Delta = -33$~nt,
$P = 1.2 \times 10^{-19}$; Supplementary Fig.~S8b).
Polyadenylation signal (PAS) strength also contributed: ancient L1
lacking any recognizable PAS showed stronger shortening
($\Delta = -35$~nt) than those with a PAS variant ($\Delta = -26$~nt;
PAS $\times$ stress $P = 0.014$).
Exon-embedded L1 within lncRNA or
pseudogene exons (29\% of total L1) were unaffected ($\Delta = -4.0$~nt,
$P = 0.33$; Supplementary Fig.~S2a).

Chromatin context further modulated the response among ancient,
non-embedded L1. L1 in active regulatory chromatin (enhancers and active
promoters; ChromHMM 15-state, HeLa-S3 E117) showed the strongest
shortening ($\Delta = -73$~nt vs.\ $-24$~nt for non-regulatory L1,
Fig.~1f; OLS stress$\times$regulatory $\beta = -36.4$,
$P = 1.1 \times 10^{-4}$), with enhancer-associated L1 most affected
($\Delta = -78$~nt, $P = 3.8 \times 10^{-7}$). This vulnerability was
additive with \msix{}, as regulatory L1 showed greater shortening across
all \msix{} quartiles (Supplementary Fig.~S3a). Within stressed
regulatory L1, the \msix{}--poly(A) association was strong (Q4 median
59~nt vs.\ 28~nt for Q1, $P = 2.9 \times 10^{-4}$; Supplementary
Fig.~S3b). At the gene level, 69 regulatory L1 host genes with reads in
both conditions showed a bidirectional response, and the \msix{}/kb of
each gene's L1 reads predicted the direction (Spearman $\rho = 0.38$,
$P = 0.001$).

The shortening was genome-wide rather than driven by a few
high-abundance loci, was independent of
concurrent $5'$ truncation at the single-molecule level ($r = 0.044$,
$P = 0.12$).

To test whether shortening requires stress granule (SG) formation,
we examined cycloheximide (CHX) co-treatment, which prevents SG
assembly by trapping ribosomes on mRNAs. CHX rescued L1 poly(A)
shortening: Median poly(A) under Ars+CHX was
93.6~nt, indistinguishable from mock (88.2~nt, $P = 0.66$), compared
with 65.9~nt under arsenite alone ($P = 0.005$ vs.\ mock, Fig.~1d). The
rescue was specific to ancient L1 ($\Delta = +29.3$~nt,
$P = 1.3 \times 10^{-4}$) and uniform across \msix{} quartiles
, placing SG formation upstream of the
\msix{}-associated poly(A) variation described below. L1 RPM was 1.51-fold
higher under Ars+CHX than arsenite alone (879 vs.\ 582).

We next assessed $5'$ decay. XRN1 knockdown
alone increased L1 proportion 1.46-fold ($P = 2.0 \times 10^{-11}$,
$\chi^2$ test), while XRN1 knockdown under arsenite had no additional
effect (1.07-fold, $P = 0.15$), indicating convergence on the same
pathway (Fig.~1e). L1 \msix{}/kb remained constant across all four
conditions ($\sim$2.5/kb). Independent ribo-depleted Illumina RNA-seq
under the same arsenite treatment (GSE278916, $n = 2$ per condition)
showed reduced RNA abundance, with 53 of 67 ancient L1 subfamilies with
significant changes showing downregulation (Supplementary Fig.~S5).

% ============================================================
% RESULTS Part 3
% ============================================================
\subsection{\msix{} density is associated with L1 poly(A) retention under stress}

We next asked whether \msix{} density is associated with the extent of
poly(A) shortening on individual L1 molecules. Although \msix{} density
itself showed only a modest decrease under arsenite (2.50 vs.\ 2.62
sites/kb, 4.4\% reduction, Mann--Whitney $P = 0.012$), the relationship
between \msix{} density and poly(A) length changed markedly. We
stratified L1 reads into \msix{}/kb quartiles and compared poly(A)
distributions (Fig.~2a). In unstressed HeLa cells, poly(A) tail length
varied modestly across quartiles (Q1 = 109 nt, Q4 = 140 nt,
$\Delta$ = 31 nt), whereas under arsenite the separation widened, with
low-\msix{} reads (Q1) showing a median poly(A) of 65 nt and
high-\msix{} reads (Q4) retaining 124 nt ($\Delta$ = 59 nt).

Under normal conditions, \msix{}/kb and poly(A) were uncorrelated
($r \approx 0$; Supplementary Fig.~S2b), whereas under arsenite a
positive association emerged (per-read Spearman $\rho = 0.201$,
$P = 3.3 \times 10^{-26}$; median $\Delta$poly(A) from $-42$~nt in Q1
to $-16$~nt in Q4; Supplementary Fig.~S6). PAS presence was an
independent protective factor: when both PAS and \msix{} were included
in the same model, each
remained significant (PAS $\times$ stress $P = 0.014$, \msix{}/kb
$\times$ stress $P = 1.4 \times 10^{-4}$). The coupling was driven by
ancient L1, where the \msix{}--poly(A) correlation increased from
$r = 0.06$ (normal) to $r = 0.18$ ($P < 10^{-19}$) under stress, while
young L1 showed no change ($r \approx 0.07$, ns; Fig.~2b). Stratifying
stressed ancient L1 by \msix{} and PAS status (CpG density averaged)
revealed a 69-nt range in median poly(A), from 132~nt for high-\msix{}
with canonical PAS to 63~nt for low-\msix{} without PAS (32\%
below 30~nt; Supplementary Fig.~S6b).

Reads with critically short poly(A) tails ($<$ 30 nt, below the PABPC1
dissociation threshold) were enriched among low-\msix{} L1 under stress
Among Q1 reads, 30.6\% fell below this threshold, compared with
15.3\% of Q4 reads (2.0-fold, Fisher's exact test,
$P = 5.9 \times 10^{-11}$).
Under normal conditions, this disparity was modest (Q1 14.6\% vs.\ Q4
10.3\%, $P = 0.036$). To estimate how much poly(A) variation is
attributable to \msix{} density, we aggregated reads by genomic L1 locus
(median \msix{}/kb and median poly(A) per locus; loci with $\geq$5
reads). The per-locus Spearman correlation was $\rho = 0.51$
($P = 1.7 \times 10^{-4}$, $n = 49$ loci), with an OLS $R^2 = 0.12$ and a
weighted least-squares $R^2 = 0.18$ (Supplementary Table~S1). The
effect was absent in unstressed loci ($R^2 = 0.003$,
$n = 140$), confirming stress dependence. Together, these data indicate
that L1 poly(A) retention under stress is associated with a hierarchy of
molecular features: PAS strength as the primary constitutive factor,
\msix{} density as a stress-specific correlate, and
chromatin context as an independent modulator.

% --- Figure 2 ---
\begin{figure}[H]
\centering
\includegraphics[width=\textwidth]{fig2.pdf}
\caption{\textbf{\msix{} density is associated with L1 poly(A) retention under stress.}
See Figure Legends for full description.}
\label{fig:m6a_polya_coupling}
\end{figure}

% ============================================================
% RESULTS Part 4
% ============================================================
\subsection{Sequence features of young L1 confer stress immunity}

Having established that \msix{} density covaries with poly(A) retention
under stress, we asked what sequence-encoded features distinguish young
L1, which are immune to shortening, from vulnerable ancient elements.
Young L1 differed from ancient elements across five axes: \msix{}/kb
(1.8-fold), full-length fraction (19.9-fold), EN domain retention
(4.0-fold), consensus span (8.8-fold), and $3'$~UTR coverage
(1.6-fold; Supplementary Fig.~S8a). Each of these features
independently conferred stress immunity when present in ancient L1:
full-length ancient elements ($\Delta = -9$~nt, ns), EN-containing
elements ($\Delta = -6.1$~nt, ns), elements with high \msix{}
($\geq$ young median; $\Delta = -12.6$~nt, ns), and elements with long
consensus span ($\Delta = +5.6$~nt, ns) were all resistant, whereas
ancient elements lacking each feature showed significant shortening
($\Delta = -34$ to $-47$~nt, all $P < 0.001$;
Supplementary Fig.~S8b). The sole exception was $3'$~UTR coverage,
which did not protect ancient elements ($\Delta = -24.4$~nt,
$P = 4.3 \times 10^{-7}$), consistent with the requirement for
upstream structural features beyond the PAS. A composite immunity
score (0--3, one point each for full-length status, EN domain
retention, and high \msix{}/kb) showed a dose-dependent relationship
with poly(A) retention under stress: score~0 median 66~nt, score~1
91~nt, score~2 132~nt, score~3+ 137~nt (Spearman $\rho = 0.169$,
$P = 5.5 \times 10^{-18}$; Supplementary Fig.~S8c).

At the motif level, the EN domain harbors the highest DRACH density in
L1 (43.2/kb, 1.38-fold the consensus average) alongside near-complete
CpG depletion (2.8/kb, obs/exp = 0.08; Supplementary Fig.~S8d),
creating a region maximally accessible to METTL3-mediated \msix{}
deposition while invisible to ZAP-mediated CpG surveillance
\citep{Takata2017,Moldovan2015}. Consistent with this, mutation
sensitivity mapping showed that stress-sensitive ancient L1 carry
significantly higher mutation rates specifically in the EN domain
. \msix{}/kb retained independent predictive value after
controlling for structural features. Among stressed ancient L1 reads
($n = 2{,}572$), a likelihood ratio test comparing OLS models with and
without \msix{}/kb (controlling for full-length status, EN domain,
consensus span, and $3'$~UTR coverage) yielded $\chi^2 = 57.9$,
$P = 2.7 \times 10^{-14}$, with each unit increase in \msix{}/kb
associated with $+6.3$~nt longer poly(A) ($\Delta R^2 = 2.2\%$). The
larger coefficient compared with the pooled OLS
($\beta = 4.2$, Supplementary Table~S2) reflects the restriction to
stressed ancient elements, where the \msix{}--poly(A) association is
strongest. These results
place \msix{} density alongside structural completeness and PAS
strength in a multi-feature framework that quantitatively predicts L1
poly(A) fate under stress.


% ============================================================
% DISCUSSION
% ============================================================
\section{Discussion}

We show that young, full-length L1 retain their poly(A) tails under
arsenite stress, and that overlapping sequence features---structural
completeness, EN domain retention, high \msix{} density, and long
consensus span---each independently confer this protection. Ancient L1
lacking these features undergo selective poly(A) shortening, and among
these vulnerable elements, \msix{} density covaries with poly(A)
retention in a dose-dependent pattern. The \msix{} marks that promote
retrotransposition \citep{Hwang2021} and mediate transcriptional
silencing \citep{Liu2021,Tian2024} thus show a third, stress-dependent
dimension: an association with differential poly(A) retention among
structurally incomplete L1.

L1 poly(A) shortening under arsenite is unexpected given that canonical
deadenylation is suppressed: CCR4-NOT is sequestered into stress
granules \citep[preprint]{Bhatt2025} and the deadenylase adaptors Tob
and Pan3 are degraded \citep{Funakoshi2015}. Non-L1 transcripts
accordingly retain their poly(A) (Fig.~1b), making L1 a notable
exception. Although
retrotransposon poly(A) tails can form protective triple-helix
structures in plant systems \citep{Li2025triplehelix}, human L1 lacks
the ENE motif required for such structures, and its poly(A) tails are
clearly vulnerable to shortening. L1, as non-translated, intronless,
AU-rich RNA, may instead engage non-canonical decay. The HUSH--NEXT
axis targets intronless transcripts for nuclear decay
\citep{Tchasovnikarova2022,Garland2022}, which would explain the
resistance of exon-embedded L1 that inherit spliced architecture.
MOV10--DCP2 mediates $5'$ decay through phase-separated granules
\citep{LiXembr2023}, a mechanism that may underlie the concurrent $5'$
truncation we observe. Cycloheximide blocks stress granule assembly and
rescues poly(A) (Fig.~1d), placing shortening downstream of SG
formation, as reported for SG-dependent $5'$ decay \citep{deMuro2024}.

L1 additionally undergoes $3'$ poly(A) shortening, a process distinct
from $5'$ decay: the two are independent at the single-molecule level,
and XRN1---which can process \msix{}-containing RNA without impediment
\citep{Athapattu2021}---degrades L1 non-selectively (L1 \msix{}/kb
unchanged across all four factorial conditions). In contrast, poly(A)
shortening correlates with \msix{} density and is strongest in
regulatory chromatin (Fig.~1f). These axes separate bulk turnover
(XRN1) from stress-inducible poly(A) shortening that is associated
with \msix{} density. Supporting this dichotomy, L1 reads experiencing
$3'$-only vulnerability carried lower \msix{}/kb than reads affected by
neither pathway (1.42 vs.\ 2.04, $P = 1.0 \times 10^{-15}$), whereas
XRN1 sensitivity was uniform across \msix{} strata. The coupling was
similar in intronic and intergenic L1 ($r \approx 0.20$ for both),
indicating that selectivity reflects transcript-intrinsic features
rather than genomic context.

How might a modification associated with decay under basal conditions
become associated with retention under stress? Arsenite relocates
YTHDF1/3 to SG peripheries and recruits YTHDF2 for structural
maintenance rather than decay \citep{FuZhuang2020,LiuSciAdv2024},
potentially allowing IGF2BP stability-promoting readers
\citep{Huang2018igf} to dominate. However, L1 may be a poor YTHDF2
substrate given its lack of translational context, and YTHDF2 phase
separation under arsenite \citep{Wang2023ythdf2} would predict
enhanced decay rather than protection---distinguishing these models
will require reader-specific perturbation. The \msix{}--poly(A)
association on L1 inverts the canonical mRNA relationship, where
\msix{} promotes cotranslational decay via CCR4-NOT recruitment
\citep{Lima2017,Absmeier2024}. Ancient L1, presumably untranslated,
escape this pathway, which may explain their longer baseline poly(A)
(122 vs.\ 107~nt) but also deprive them of ribosome-mediated
protection under stress.

Young L1 stress immunity reflects overlapping sequence-encoded
protections rather than a single determinant
(Supplementary Fig.~S8). As domains are lost through $5'$ truncation,
these protections are progressively stripped, and \msix{} emerges as a
late-stage correlate whose association with poly(A) retention becomes
detectable only in structurally incomplete elements that lack other
protections. This framing aligns with the \msix{} surveillance model,
which proposes that \msix{} preferentially marks intronless transcripts
as quality-control targets \citep{Ke2026}. Under basal conditions, such
marking may facilitate decay; under stress, when SG sequestration
disables canonical decay, the same marks may be co-opted for
differential poly(A) retention.

These observations have implications for L1 regulatory biology.
Supporting a functional role for the \msix{} enrichment we report,
\msix{} on ancient L1PA elements directly regulates
EP300/KAP1 binding and LTR trans-silencing in human embryonic stem cells
\citep{Zhu2025L1PA}. The selective poly(A) shortening we observe under
stress could therefore alter the composition of the regulatory L1 RNA
pool, although demonstrating downstream chromatin consequences will
require perturbation of \msix{} writers specifically on L1.

The stress resilience of young L1 may help explain a long-observed
phenomenon: L1 RNA accumulation under cellular stress and senescence
\citep{DeCecco2019,Gorbunova2021}. Prior studies have attributed this
accumulation to transcriptional derepression---DNA demethylation and
chromatin relaxation unlock L1 promoters
\citep{Gorbunova2021}. Our
data suggest that post-transcriptional stability is a complementary
mechanism. Young, full-length L1 carry overlapping protections
(Supplementary Fig.~S8a,b) that shield their poly(A) tails from
stress-induced shortening, allowing newly transcribed L1 RNA to
persist. In contrast, ancient L1 fragments lacking these features
undergo selective poly(A) shortening and decay. The net effect---young
L1 RNA accumulation alongside ancient L1 RNA depletion---may reshape
the cellular L1 RNA pool under stress, favoring retrotransposition-competent
over regulatory transcripts. These features are sequence-encoded and
cross-cell-line consistent (CV~=~0.112 for \msix{}/kb), offering a
quantitative basis for predicting which L1 elements are most
susceptible to stress-induced poly(A) shortening.

Several limitations constrain interpretation. The \msix{}--poly(A)
association, while robust and dose-dependent, remains correlational.
The per-read OLS model explains a modest fraction of poly(A) variance
($R^2 = 0.031$ for stressed reads; $R^2 = 0.039$ pooled,
Supplementary Table~S2), as expected given single-molecule noise and
the multiple determinants of poly(A) length. Per-locus aggregation
raises the explained variance to $R^2 = 0.12$--0.18
(Supplementary Table~S1), but rests on only 49 loci with $\geq$5 reads
(bootstrap 95\% CI for $\rho$: 0.28--0.70, 10{,}000 resamples). The association is
significant across all five probability thresholds tested
(50\% to 90\%, all $P < 2 \times 10^{-4}$; Supplementary
Table~S3), arguing against threshold-dependent circularity.
Perturbation of \msix{} writers under stress would be needed to
establish causality. We note that analysis of a published METTL3
knockout dataset (HEK293T; PRJEB40872) showed no significant reduction
in L1 \msix{}/kb (fold change = 0.97, $P > 0.05$), though non-L1
control transcripts showed a modest decrease (0.91-fold, $P = 0.035$).
This null result may reflect cell-type differences, residual METTL3
activity, or limitations of modification calling on repetitive
sequences, but it underscores the need for L1-targeted writer
perturbation under stress conditions.

DRS captures only surviving polyadenylated RNA and
cannot measure decay rates directly. The $3'$ bias of DRS (90\% of reads
within 1~kb of the $3'$ end) limits $5'$ UTR coverage; the median L1 body
overlap per read is 0.35, so read-level \msix{}/kb includes flanking
modifications, though restricting to high-overlap reads preserves the
dose-response pattern. ChromHMM states from E117 (HeLa-S3) were used as a
proxy for HeLa chromatin; subline-specific epigenetic divergence cannot
be excluded. All stress-related findings derive from a single cell line
(HeLa) treated with arsenite for 60~minutes. Baseline \msix{} enrichment
is consistent across nine cell lines (CV~=~0.112), but generalizability
to other cell types and stressors remains to be established. MAFIA was
trained on non-repetitive transcripts \citep{Pratanwanich2021,Hendra2022},
and nanopore modification callers show variable performance across
sequence contexts \citep{Deng2024bench}; benchmarking specifically on
repetitive elements remains a priority. However, orthogonal RNA004/Dorado
validation reproduced the L1 enrichment with 93\% DRACH motif specificity
(Supplementary Fig.~S6), supporting modification call reliability. Time-course and
multi-cell-line stress experiments are important next steps. For the
XRN1~KD experiments, statistical comparisons are based on read-level
variation within single biological replicates per condition due to a
sequencing-platform batch effect (see Methods).

% ============================================================
% MATERIALS AND METHODS
% ============================================================
\section{Materials and Methods}

\subsection*{Data sources}

We reanalysed publicly available Oxford Nanopore direct RNA sequencing
(DRS) data from three sources.
\textbf{(i)}~Baseline transcriptome data for nine human cell lines
(A549, H9, HCT116, HeLa, HepG2, HEYA8, K562, MCF7, and SH-SY5Y; 23
libraries total) were obtained from the SG-NEx project
\citep{ChenSGNEx2025}. All libraries used the ONT SQK-RNA002 kit
(MinION, FLO-MIN106, R9.4.1 chemistry) with at least two biological
replicates per cell line.
\textbf{(ii)}~HeLa arsenite-stress data (0.5~mM NaAsO$_2$, 60~min; three
replicates), cycloheximide co-treatment data, and XRN1 knockdown data
were obtained from PRJNA842344 \citep{deMuro2024}. For XRN1~KD/CHX,
replicate~1 (PromethION) and replicate~2 (MinION) showed a 2.94-fold
platform batch effect; only MinION replicate~2 was used for all four
factorial conditions (mock, XRN1~KD, arsenite, arsenite~+~XRN1~KD), so
comparisons are based on read-level variation within single
replicates. MCF7-EV data were from the same SG-NEx panel.
\textbf{(iii)}~For RNA004 validation, HeLa DRS data were basecalled with
Dorado v0.8.0 (ONT; model
\texttt{rna004\_130bps\_sup@v5.2.0\_inosine\_m6A\_2OmeA@v1}), which
detects \msix{} at all adenosine positions regardless of motif context.
Per-read \msix{}/kb was computed from the BAM-encoded per-base
modification probabilities (threshold $\geq 0.80$) to compare with
MAFIA-derived values. Raw signal data were obtained as FAST5 or POD5 files.

\subsection*{Basecalling and alignment}

Raw FAST5 files were basecalled using Guppy v6.0.0 (ONT) with the
high-accuracy RNA model (\texttt{rna\_r9.4.1\_70bps\_hac}). Basecalled
reads were aligned to GRCh38/hg38 using minimap2 v2.28
\citep{Li2018minimap2} with splice-aware settings and a GENCODE v38
junction BED file (\texttt{--junc-bed}). Secondary alignments were
discarded and BAM files were sorted and indexed with samtools v1.23.

\subsection*{Two-stage L1 identification filter}

L1 reads were identified through a two-stage filter designed to
eliminate host gene read-through contamination.

\emph{Stage 1 (coarse filter).} Aligned reads were intersected with
RepeatMasker L1 annotations (UCSC rmsk track, hg38) using bedtools
v2.31.0. Reads with $\geq$10\% overlap to annotated L1 elements were
retained.

\emph{Stage 2 (fine filter).} Stage 1 reads were filtered by four
criteria: (i) exclusion of reads with splice junctions (CIGAR `N'
operations $>$ 50 bp), (ii) selection of the best alignment per read,
(iii) exclusion of reads with $\geq$100 bp overlap to protein-coding
exons (GENCODE v38), and (iv) strand concordance with the annotated L1
element. Reads passing both stages were classified as non-embedded L1
(``PASS'' in the analysis pipeline). The 100~bp exon overlap threshold
was empirically validated: 92.1\% of non-embedded L1 reads have zero exon
overlap, and the biological signal (arsenite poly(A) shortening) is
indistinguishable between zero-overlap and near-threshold reads
. Reads failing Stage~2 criterion (iii) due to
overlap with long noncoding RNA or pseudogene exons (but not
protein-coding exons) were classified as exon-embedded L1.

Non-L1 reads from the same libraries were randomly subsampled to form
the comparison group and processed through identical nanopolish and
MAFIA pipelines. Because DRS captures poly(A)$^+$ RNA, these non-L1 reads
are dominated by highly expressed protein-coding transcripts: 93.9\%
protein-coding mRNA, 3.2\% MT-rRNA, and 1.7\% lncRNA ($n = 40{,}659$
reads across 6,512 genes). The top contributors include \emph{MT-RNR2},
\emph{EEF1A1}, \emph{FTL}, and ribosomal protein mRNAs. This composition
differs from L1 reads in sequence context, expression level, and
transcript architecture. To address potential bias from
expression-level differences in \msix{}/kb, we performed
expression-matched analyses using gene-weighted and
quintile-stratified controls (see Methodological considerations). These
reduced the L1/non-L1 \msix{}/kb ratio but did not alter the principal
conclusions.

As an additional control for non-coding RNA identity, annotated lncRNAs
(GENCODE v38 biotype ``lncRNA'') were independently identified from the
same libraries and processed through identical nanopolish poly(A) and
MAFIA pipelines (6 libraries $\times$ $\sim$3{,}000--4{,}800 MAFIA reads).
Read-length-matched analyses confirmed that poly(A) and \msix{}
differences were not attributable to transcript length
(Supplementary Fig.~S7).

\subsection*{L1 subfamily and age classification}

L1 subfamily assignments were derived from RepeatMasker annotations.
Elements were classified as ``young'' (L1HS, L1PA1, L1PA2, L1PA3) or
``ancient'' (all other L1 subfamilies including L1MC, L1ME, L1M families).
Consensus positions within the full-length L1 reference were obtained
from RepeatMasker \texttt{repStart}, \texttt{repEnd}, and
\texttt{repLeft} fields (UCSC rmsk track). Functional domains were
defined on the L1HS consensus (6{,}064~bp): $5'$~UTR (1--908),
ORF1 (909--1{,}990), ORF2 (1{,}991--5{,}817), and $3'$~UTR
(5{,}818--6{,}064). Each genomic L1 element was assigned a primary
domain based on the largest consensus overlap, and elements spanning
all four domains were classified as ``full-length''. For the ORF2
sub-domain analysis, the EN domain was defined as consensus positions
1{,}991--2{,}708, the RT domain as 2{,}709--4{,}149, and the C-rich
domain as 4{,}150--5{,}817.

\subsection*{Mutation sensitivity mapping}

To identify L1 consensus positions where mutations correlate with stress
vulnerability, we extracted the genomic sequences of 508 ancient L1 loci
(top and bottom poly(A) quartiles under arsenite) using
bedtools~getfasta (strand-aware). Sequences were aligned to the L1HS
consensus (GenBank L19088.1) with minimap2 (\texttt{-{}-eqx -{}-MD}).
Per-position match/mismatch counts were tallied separately for
stress-sensitive and stress-resistant loci. Mutation rate differences
were assessed in 100-bp sliding windows (50-bp step) using Fisher's
exact test with Bonferroni correction. Only windows with $\geq$100
aligned bases in both groups were included.

\subsection*{Poly(A) tail length estimation}

Poly(A) tail lengths were estimated from raw nanopore signal using
nanopolish v0.14.0 \citep{Simpson2017nanopolish} with the
\texttt{polya} module, which fits a piecewise model to the poly(A)
signal segment. Only reads with a ``PASS'' quality flag from nanopolish
were retained.

\subsection*{\msix{} modification calling}

\msix{} modifications were called from aligned BAM files using
psi-co-mAFiA (v0.0.1; \citealp{Hendra2022}), which
assigns per-base \msix{} modification probabilities (ChEBI:21891)
encoded as MM/ML tags in the BAM output.

Modification calls were parsed from MM/ML BAM tags, matching both
forward-strand (N$+$\allowbreak 21891) and reverse-strand
(N$-$\allowbreak 21891) tags to capture modifications on all mapped reads.
A probability threshold of $\geq 0.80$ (ML $\geq$ 204/255) was used to
classify sites as modification-positive for all analyses.

\subsection*{\msix{} density and decomposition analysis}

\msix{} density was computed as the number of modification-positive DRACH
sites per kilobase of aligned read length. To determine whether L1
\msix{} density exceeds what would be expected from elevated DRACH motif
frequency in AT-rich L1 sequences, we counted DRACH motif positions
(D = A/G/T, R = A/G, H = A/C/T) within each read's aligned reference
sequence. Reference sequences were extracted in a CIGAR-aware manner,
walking only through M (match), = (sequence match), and X (mismatch)
CIGAR operations to exclude intronic regions in spliced non-L1 reads.
This CIGAR-aware extraction was critical for non-L1 transcripts, where
using \texttt{reference\_end} $-$ \texttt{reference\_start} (which
includes intronic spans) inflated motif counts by approximately 4-fold.

\subsection*{Statistical analyses}

All statistical analyses were performed in Python 3.10 using scipy,
statsmodels, and pandas. Two-group comparisons used the Mann--Whitney
$U$ test (poly(A) tail lengths, \msix{}/kb) or Fisher's exact test
(\msix{} site counts). Multi-group comparisons used the
Kruskal--Wallis test. Effect sizes for modification enrichment were
computed as fold change of \msix{}/kb. The relationship between \msix{} density and
poly(A) tail length was modeled using ordinary least-squares (OLS)
regression with the model:

\begin{equation}
\text{poly(A)} = \beta_0 + \beta_1 \cdot \text{\msix{}/kb} + \beta_2 \cdot \text{stress} + \beta_3 \cdot (\text{\msix{}/kb} \times \text{stress}) + \epsilon
\end{equation}

\noindent where stress is a binary indicator for the arsenite condition.
The full model included additional covariates for read length (z-scored),
L1 age (young vs.\ ancient), and their interactions
(Supplementary Table~S2).
The interaction coefficient $\beta_3$ tests whether the \msix{}--poly(A)
relationship differs between stressed and unstressed conditions. Pearson
and Spearman correlation coefficients were computed for stratified
subgroups (L1 age, genomic context, stress condition).

To reduce single-molecule measurement noise, we performed per-locus
aggregation by grouping reads that mapped to the same genomic L1
insertion (defined by chromosome, start, end, and strand). For each
locus, median \msix{}/kb and median poly(A) length were computed.
Per-locus regression was performed using both ordinary least-squares
and weighted least-squares (weight = number of reads per locus) on
loci with a minimum read depth of 3 or 5 reads. Multiple testing
correction was not applied to the primary analyses, as each test
addressed a distinct, pre-specified hypothesis. Where exploratory
analyses involved multiple comparisons, this is noted in the text.

\subsection*{Illumina RNA-seq validation}

To orthogonally validate L1 expression changes under arsenite,
we analysed publicly available ribo-depleted Illumina RNA-seq data from
HeLa cells (GSE278916, untreated $n=2$, arsenite $n=2$). Reads were
aligned to the human genome (hg38, GENCODE v38) using STAR v2.7.11a with
multi-mapper settings compatible with TE quantification
(\texttt{--outFilterMultimapNmax 100}). L1 subfamily-level expression
was quantified with TEtranscripts v2.2.3 using the EM algorithm for
multi-mapper resolution and the RepeatMasker TE annotation
(hg38\_rmsk\_TE.gtf). Differential expression was assessed with
DESeq2 (integrated within TEtranscripts). As a complementary approach,
element-level L1 counts were obtained with featureCounts (Subread v2.0.6)
in fractional multi-mapper mode and normalized using DESeq2-style
median-of-ratios gene-level size factors.

\subsection*{Gene ontology and genomic context analysis}

Host genes harboring intronic L1 were identified by intersecting L1
genomic coordinates with GENCODE v38 gene annotations. Gene Ontology
enrichment was performed using Enrichr. PAS motifs were scanned within
50 bp upstream of each L1 element's $3'$ end in the reference genome,
searching for the canonical AATAAA and 11 variant PAS hexamers.

\subsection*{ChromHMM chromatin state annotation}

Chromatin state annotations were obtained from the Roadmap
Epigenomics 15-state ChromHMM model \citep{Roadmap2015}. For HeLa
analyses, the E117 (HeLa-S3) segmentation BED file was used. Each L1
element was assigned to a chromatin state by intersecting its genomic
midpoint with the ChromHMM state annotation using bedtools intersect.
States were grouped into functional categories: ``regulatory'' (active
enhancers [states 6--7] and active promoters [states 1--3]),
``transcribed'' (states 4--5, 9--11), and ``quiescent/repressed''
(states 12--15).

% ============================================================
% DATA AVAILABILITY
% ============================================================
\section{Data Availability}

This study reanalysed publicly available datasets; no new sequencing
data were generated. Baseline DRS data for nine human cell lines were
obtained from the SG-NEx project (ENA: ERP125816). Arsenite-stress,
cycloheximide, and XRN1~KD data were from PRJNA842344
\citep{deMuro2024}. PUS enzyme knockdown data were from PRJNA1220613
(BE(2)-C) and PRJNA1092333 (SH-SY5Y). Analysis code and
figure-generation scripts are available at \url{https://github.com/junsoopablo/L1-m6A-polyA}.

% ============================================================
% ACKNOWLEDGMENTS
% ============================================================
\section*{Acknowledgments}

We thank Hyeshik Chang for computational server resources and helpful
comments on the manuscript. This work was not supported by external
funding.

\subsection*{Author Contributions}
J.K.\ conceived and designed the study, performed all computational
analyses, and wrote the manuscript. S.L.\ acquired and curated the
public datasets. S.H.\ contributed to data interpretation. All authors
reviewed the manuscript.

% ============================================================
% REFERENCES
% ============================================================
\bibliography{references}

% ============================================================
% FIGURE LEGENDS
% ============================================================
\clearpage
\section*{Figure Legends}

\subsection*{Figure 1. L1 \msix{} enrichment and arsenite-selective poly(A) shortening.}

\textbf{(a)} ECDF of per-read \msix{}/kb for young L1 (L1HS,
L1PA1--3), ancient L1, and non-L1 transcripts. Dotted vertical lines
mark group medians. L1 curves are right-shifted, indicating higher
\msix{} density (1.79-fold enrichment).
%
\textbf{(b)} ECDF of poly(A) tail length for L1 and non-L1 transcripts
in unstressed HeLa versus arsenite-treated HeLa. The L1 curve shifts
leftward under arsenite ($\Delta = -31$~nt, $P = 2.6 \times 10^{-17}$)
while non-L1 curves remain unchanged.
%
\textbf{(c)} Poly(A) distributions for ancient and young L1 under
normal and arsenite conditions. Ancient L1 show shortening
($\Delta = -33$~nt, $P = 1.7 \times 10^{-18}$) while young L1 are
unaffected ($\Delta = -0.1$~nt, ns).
%
\textbf{(d)} Cycloheximide (CHX) rescue. Ars+CHX poly(A) is
indistinguishable from mock ($P = 0.66$), placing shortening downstream
of stress granule formation.
%
\textbf{(e)} XRN1 knockdown increases L1 proportion 1.46-fold under
normal conditions ($P = 2.0 \times 10^{-11}$) but has no additional
effect under arsenite (1.07-fold, $P = 0.15$), indicating pathway
convergence. Error bars: 95\% CI.
%
\textbf{(f)} Chromatin state modulates arsenite vulnerability. L1 in
regulatory chromatin show the strongest shortening ($\Delta = -73$~nt),
while transcribed and quiescent regions show moderate effects.
ChromHMM 15-state, HeLa-S3 E117.

\subsection*{Figure 2. \msix{} density is associated with L1 poly(A) retention under stress.}
\textbf{(a)} ECDF of poly(A) length stratified by \msix{}/kb quartile.
Under stress, low-\msix{} reads (Q1) shift leftward while high-\msix{}
reads (Q4) retain poly(A). Dotted line: 30-nt PABP dissociation
threshold (Q1/Q4 ratio 2.0-fold, $P = 5.9 \times 10^{-11}$).
%
\textbf{(b)} Pearson $r$ (\msix{}/kb vs.\ poly(A)) under normal versus
stress for young L1, ancient L1, and all L1. The correlation emerges
under stress in ancient L1 ($r = 0.06 \to 0.18$, $P < 10^{-19}$),
while young L1 shows no change.


\end{document}
